% Figure: two pipes in parallel between a common upstream node A and a common
% downstream node B. Both branches carry the same head loss (the pressure drop
% from A to B is shared), and the total flow splits between them in inverse
% proportion to their resistances. Drawn as a schematic loop with the flow
% entering at A, splitting into branch 1 (larger bore, lower resistance) and
% branch 2 (smaller bore, higher resistance), and recombining at B.
\begin{figure}[htbp]
\centering
\begin{tikzpicture}[>={Latex}, thick]
% Nodes A and B.
\node[circle, draw=engblack, fill=enggray!12, minimum size=7mm] (A) at (0,0) {$A$};
\node[circle, draw=engblack, fill=enggray!12, minimum size=7mm] (B) at (9,0) {$B$};
% Inlet and outlet stubs.
\draw[->, engblue] (-1.7,0) -- (A);
\node[font=\scriptsize, engblue, anchor=south] at (-1.2,0.05) {$Q$};
\draw[->, engblue] (B) -- (10.7,0);
\node[font=\scriptsize, engblue, anchor=south] at (10.2,0.05) {$Q$};
% Branch 1 (upper, larger bore): thicker line.
\draw[engred, line width=1.6pt] (A) .. controls (3,1.9) and (6,1.9) .. (B);
\node[font=\scriptsize, engred] at (4.5,2.15)
  {branch 1: $D_1$, $L_1$, $f_1$ \quad $Q_1$};
% Branch 2 (lower, smaller bore): thinner line.
\draw[engamber, line width=0.9pt] (A) .. controls (3,-1.9) and (6,-1.9) .. (B);
\node[font=\scriptsize, engamber] at (4.5,-2.2)
  {branch 2: $D_2$, $L_2$, $f_2$ \quad $Q_2$};
% Head-loss annotation: same drop across both branches.
\node[font=\scriptsize, enggray, anchor=center] at (4.5,0)
  {$h_{f,1} = h_{f,2}$, \quad $Q = Q_1 + Q_2$};
\end{tikzpicture}
\caption[Two pipes in parallel]{Two pipes in parallel between a common
upstream node $A$ and downstream node $B$. The head loss across the pair is the
same along either branch, $h_{f,1} = h_{f,2}$, because both branches connect the
same two nodes, and the total flow is the sum of the branch flows, $Q = Q_1 +
Q_2$. These two statements, equal head loss and conserved flow, are the
parallel-network analogue of the series rule that adds the drops along one line;
they close the split of the flow between the branches once each branch
resistance is known. Worked example~34 carries the numbers for a larger bore
and a smaller bore sharing a duty flow.}
\label{fig:vol03ch12:parallel-pipes}
\end{figure}
